\documentclass[12pt]{article}

\usepackage[margin=1in]{geometry}
\usepackage{amsmath, amssymb}
\usepackage{booktabs}
\usepackage{setspace}
\usepackage{enumitem}

\setstretch{1.2}
\setlength{\parindent}{0pt}

\title{\textbf{Solutions  In-Class Group Exercise}\\
Labour Markets and the Future of Work Exercise}
\author{ECON 204 -- Contemporary Economic Issues}
\date{}

\begin{document}

\maketitle

\section*{Part I: Structural Change in Employment}

\textbf{Question 1}

Manufacturing experienced the largest decline in employment share.

\textbf{Question 2}

Manufacturing fell from $25\%$ to $10\%$.

\[
25\% - 10\% = 15 \text{ percentage points}
\]

\textbf{Answer:} 15 percentage points.

\textbf{Question 3}

The table illustrates \textbf{structural transformation}. Over time, advanced economies tend to shift:

\begin{itemize}
\item from agriculture,
\item to manufacturing,
\item and eventually toward services.
\end{itemize}

This pattern occurs because productivity growth reduces labour demand in goods-producing sectors while rising incomes increase demand for services.

\textbf{Question 4}

Two key explanations for the decline in manufacturing employment are:

\begin{enumerate}
\item \textbf{Automation and technological change}. Modern machinery and robotics allow firms to produce more output with fewer workers.
\item \textbf{Globalization and offshoring}. Labour-intensive production stages have moved to countries with lower labour costs.
\end{enumerate}

Additional factors may include rising productivity and changing consumer demand toward services.

\textbf{Question 5}

Service-sector employment has become dominant because:

\begin{itemize}
\item rising incomes increase demand for services,
\item productivity growth reduces labour demand in agriculture and manufacturing,
\item many service occupations involve interpersonal or non-routine tasks,
\item aging populations increase demand for health and care services.
\end{itemize}

\section*{Part II: Automation, Tasks, and Occupations}

\textbf{Question 6}

The occupations most vulnerable to automation are those dominated by routine tasks.

\begin{itemize}
\item Data Entry Clerk --- highly routine and easily automated
\item Taxi Driver --- vulnerable due to autonomous vehicle technologies
\end{itemize}

Moderately vulnerable occupations include accountants because some tasks can be automated.

Less vulnerable occupations include nurses and software developers, which require interpersonal interaction and complex problem solving.

\textbf{Question 7}

Routine tasks are easier to automate because they follow predictable and rule-based procedures. Computers and algorithms perform well when tasks involve repetitive steps that can be coded into software.

Examples include:

\begin{itemize}
\item bookkeeping
\item data processing
\item assembly-line tasks
\end{itemize}

\textbf{Question 8}

Tasks requiring interpersonal skills, creativity, and judgment are harder to automate because they involve:

\begin{itemize}
\item empathy and human interaction
\item adapting to unexpected situations
\item complex decision making
\item creativity and innovation
\end{itemize}

Examples include teaching, nursing, and management.

\textbf{Question 9}

Automation often replaces specific tasks rather than entire occupations. Most jobs involve multiple tasks, some routine and some non-routine. Technology may automate the routine parts while workers continue performing tasks requiring judgment, communication, and creativity.

\section*{Part III: Wage Inequality and Skill Demand}

\textbf{Question 10}

Wage gap in 2000:

High-skill wage: \$30/hour

Low-skill wage: \$18/hour

\[
30 - 18 = 12
\]

\textbf{Answer:} \$12 per hour.

\textbf{Question 11}

Wage gap in 2025:

High-skill wage: \$52/hour

Low-skill wage: \$22/hour

\[
52 - 22 = 30
\]

\textbf{Answer:} \$30 per hour.

\textbf{Question 12}

The wage gap increased from \$12 to \$30 per hour. This indicates that wage inequality between high-skill and low-skill workers has increased.

\textbf{Question 13}

\textbf{High-skill workers}

Initial wage: \$30

Final wage: \$52

\[
52 - 30 = 22
\]

\[
\frac{22}{30} \times 100 = 73.33\%
\]

\textbf{Increase:} approximately $73.3\%$.

\textbf{Low-skill workers}

Initial wage: \$18

Final wage: \$22

\[
22 - 18 = 4
\]

\[
\frac{4}{18} \times 100 = 22.22\%
\]

\textbf{Increase:} approximately $22.2\%$.

\textbf{Question 14}

Two explanations for the widening wage gap include:

\begin{itemize}
\item \textbf{Skill-biased technological change}: technology complements skilled workers and raises their productivity.
\item \textbf{Rising demand for education and specialized skills}: modern economies reward workers with advanced technical and analytical abilities.
\end{itemize}

Other explanations may include globalization, declining unionization, and labour market polarization.

\section*{Part IV: Canadian Labour Market Data}

\textbf{Question 15}

Youth workers (ages 15--24) had the highest unemployment rate in both years.

\textbf{Question 16}

Youth unemployment increased from $11.0\%$ to $13.5\%$.

\[
13.5 - 11.0 = 2.5
\]

\textbf{Answer:} 2.5 percentage points.

\textbf{Question 17}

Young workers are more vulnerable to labour market disruptions because:

\begin{itemize}
\item they have less work experience,
\item they often hold temporary or entry-level jobs,
\item they are concentrated in sectors sensitive to economic cycles,
\item they have less bargaining power.
\end{itemize}

\textbf{Question 18}

Technological change affects age groups differently.

Younger workers may adapt more easily because they possess stronger digital skills and greater flexibility.

Older workers may face higher retraining costs and may find occupational transitions more difficult late in their careers.

\section*{Part V: The Gig Economy}

\textbf{Question 19}

Advantages of gig work include:

\begin{itemize}
\item flexible working hours,
\item easy entry into the labour market,
\item opportunities for supplemental income.
\end{itemize}

\textbf{Question 20}

Disadvantages include:

\begin{itemize}
\item lack of job security,
\item absence of employer-provided benefits,
\item income volatility.
\end{itemize}

\textbf{Question 21}

A strong answer recognizes trade-offs.

Treating gig workers as employees would provide labour protections such as minimum wages and benefits. However, it could reduce flexibility and increase costs for firms.

A hybrid approach may combine flexibility with basic protections.

\textbf{Question 22}

Possible policy responses include:

\begin{itemize}
\item portable benefits systems,
\item minimum earnings standards,
\item improved worker classification rules,
\item collective bargaining rights,
\item greater transparency in platform algorithms.
\end{itemize}

\section*{Part VI: Artificial Intelligence and Work}

\textbf{Question 23}

AI complements labour when it increases worker productivity instead of replacing workers. Technology assists workers in completing tasks more efficiently.

Examples include AI diagnostic tools assisting doctors or AI software helping programmers write code.

\textbf{Question 24}

Example: doctors using AI to analyze medical scans and improve diagnostic accuracy.

\textbf{Question 25}

Example: data entry clerks, where repetitive tasks can be automated using software and artificial intelligence.

\textbf{Question 26}

Skills likely to become more valuable include:

\begin{itemize}
\item analytical thinking,
\item problem solving,
\item communication skills,
\item creativity,
\item adaptability,
\item digital literacy.
\end{itemize}

These skills complement new technologies and are harder to automate.

\section*{Part VII: Policy Debate}

\textbf{Question 27}

Policy effects depend on the chosen option.

For example, retraining programs can help displaced workers acquire new skills and transition to growing sectors.

\textbf{Question 28}

Firms may benefit from a more skilled labour force, but they may also face higher costs depending on the policy.

\textbf{Question 29}

Possible unintended consequences include:

\begin{itemize}
\item high fiscal costs,
\item mismatch between training and labour demand,
\item reduced hiring if labour regulations become too restrictive.
\end{itemize}

\textbf{Question 30}

A strong recommendation clearly states a position and explains the economic reasoning behind it while acknowledging trade-offs.

\section*{Key Takeaways}

Students should understand that:

\begin{itemize}
\item labour markets evolve through structural transformation,
\item automation often replaces tasks rather than entire jobs,
\item technological change can increase wage inequality,
\item gig work introduces flexibility but also insecurity,
\item artificial intelligence may both complement and substitute labour,
\item public policy plays an important role in shaping future labour market outcomes.
\end{itemize}

\end{document}